\chapter{The Pythonic Functional Toolbelt}

In the preceding chapters, we explored the foundations of functional programming, understanding pure functions, higher-order functions, and lazy evaluation. We have contrasted the strict, multi-paradigm approach of Python with the purely functional nature of Haskell. In this chapter, we bridge the gap. We will examine the functional utilities built into Python, specifically the \texttt{functools} module, and powerful declarative features like comprehensions. We will draw direct parallels to their native Haskell counterparts.

\section{The \texttt{functools} Module}

The Python standard library provides the \texttt{functools} module, offering higher-order functions that act on or return other functions. These functional tools emulate the natural capabilities found in purely functional languages.

\subsection{Partial Application with \texttt{partial}}

In mathematics and functional programming, partial application refers to the process of fixing a number of arguments to a function, thereby producing another function of smaller arity. We saw how Haskell natively curries all functions. In Python, we can achieve partial application explicitly using \texttt{functools.partial}.

Let us examine how to fix an exponent in a power function to create specialized functions like \texttt{square} and \texttt{cube}:

\lstinputlisting[language=Python, caption={Partial application in Python}]{../code/python/chapter8/functools_examples.py}

In Haskell, this process is natural. A function taking two arguments actually takes one argument and returns a function taking the second. We can partially apply arguments by simply providing one parameter, or by using operator sections:

\lstinputlisting[language=Haskell, caption={Partial application with operator sections in Haskell}]{../code/haskel/chapter8/partial_app.hs}

\subsection{Memoization and \texttt{lru\_cache}}

A profound consequence of referential transparency is the ability to securely cache function results. Since a pure function yields the same output for identical inputs and produces no side effects, there is no need to re-evaluate it if the inputs have been seen before.

Python provides \texttt{functools.lru\_cache} to automatically memoize pure functions. This is highly effective for computationally expensive recursive functions, such as the naive Fibonacci sequence generator. 

Because Haskell employs lazy evaluation (call-by-need), identical expressions can be shared and evaluated only once. Furthermore, we can define infinite data structures that inherently act as dynamic caches without the need for an explicit decorator.

\section{Comprehensions: Declarative Data Transformation}

List, dictionary, and set comprehensions are declarative constructs that allow us to create new collections based on existing iterables. 

\subsection{The Mathematics of Comprehensions}

Comprehensions are heavily inspired by mathematical set builder notation. For instance, the set of all squares of even numbers in a domain $D$ can be expressed mathematically as:
\[
S = \{ x^2 \mid x \in D, x \equiv 0 \pmod{2} \}
\]
In Python, this maps directly to a list comprehension, elegantly combining the operational concepts of \texttt{map()} and \texttt{filter()} into a highly readable syntax.

\lstinputlisting[language=Python, caption={List, set, and dictionary comprehensions in Python}]{../code/python/chapter8/comprehensions.py}

This concise declarative notation is inherited almost directly from Haskell's list comprehensions:

\lstinputlisting[language=Haskell, caption={List comprehensions in Haskell}]{../code/haskel/chapter8/list_comp.hs}

\section{Generator Expressions and Laziness}

Finally, Python provides generator expressions. Unlike list comprehensions that materialize the entire collection in memory, generator expressions yield elements one at a time. This represents a form of lazy evaluation in Python, very akin to Haskell's inherent laziness, allowing operations on conceptually infinite sequences without memory exhaustion.
